\section{Discussion}
\label{sec:discussion}

Cette section discute les résultats expérimentaux et présente les découvertes concernant le comportement des solveurs Chuffed et Gecode.

\subsection{Performance des solveurs}

Les résultats montrent des différences importantes entre Chuffed et Gecode. Sur les instances medium\_02 et medium\_03, Chuffed prouve rapidement (en 2.6s et 4.0s) que ces instances sont insatisfiables (aucune solution n'existe), tandis que Gecode timeout après 300 secondes sans pouvoir conclure. Cela suggère que Chuffed est plus efficace pour détecter rapidement les conflits de contraintes.

Sur les instances résolubles (small\_01 et small\_02), Gecode trouve des solutions de meilleure qualité : objectif de 0 vs 514 pour small\_01, et 3336 vs 4564 pour small\_02. Cela suggère que Gecode explore l'espace de recherche de manière plus systématique pour trouver des solutions optimales, tandis que Chuffed se contente de solutions réalisables plus rapidement.

\subsection{Facteurs de difficulté}

Les résultats montrent un pattern clair : les instances small (8-12 opérations) sont résolubles, tandis que les instances medium et large sont principalement insatisfiables. Plusieurs facteurs expliquent cette observation :

\begin{itemize}
    \item Les fenêtres temporelles peuvent être trop serrées pour permettre une planification satisfaisante avec beaucoup d'opérations
    \item Le ratio ressources/demande peut être insuffisant pour les instances plus grandes
    \item Les contraintes de précédence et de compatibilité voie-équipement peuvent créer des conflits impossibles à résoudre
\end{itemize}

\subsection{Limites et perspectives}

Les limites de cette étude incluent le nombre limité d'instances (8) et seulement deux solveurs testés. Les perspectives futures incluent tester d'autres solveurs (OR-Tools, CP-SAT, etc.) et générer plus d'instances avec différentes caractéristiques pour mieux comprendre les facteurs de difficulté.
